% Source: https://artofproblemsolving.com/community/c3h2325389_mathematics
% Stable identity: aops-topic-2325389-q01
By definition, $a\equiv b\pmod 7$ means that $7$ divides $a-b$.

\begin{itemize}
\item $15-1=14=2\cdot7$, so $15\equiv1\pmod 7$ is true.
\item $15-2=13$ is not a multiple of $7$, so $15\equiv2\pmod7$ is \textbf{false}. This is the error.
\end{itemize}
The second step itself is valid: congruence is transitive, so from two true statements $15\equiv a$ and $15\equiv b\pmod 7$ one may conclude $a\equiv b\pmod7$. Here the false premise $15\equiv2$ leads to the false conclusion $1\equiv2\pmod7$ ($7$ does not divide $1$).

Since $15=2\cdot7+1$, the remainder is $\boxed{1}$.
